% File   : LaTeX template
% Author : cyma <cyma.git@protonmail.ch>
% Source : <https://github.com/cyma/debian>
% License : MIT Copyright (c) 2020 cyma


%-----------------------------------
% PACKAGES AND CONFIGURATIONS
%-----------------------------------

\documentclass[a4paper,twoside,11pt]{article}
\usepackage[utf8]{inputenc} %inputting international characters
\usepackage[T1]{fontenc} % output font encoding
\usepackage[spanish]{babel} %change to your language
\usepackage{xcolor} % color package
\usepackage{fancyhdr} %  page margings and sizes, headers and footers,and the proper placement of figures and tables
\usepackage{hyperref} % create hyperlinks
\usepackage{extramarks} % combine with fancyhdr to give extra customization
\usepackage{geometry} %  customize page layout, implementing auto-centering and auto-balancing mechanisms
\usepackage{graphicx} % key-value interface for optional arguments to the \includegraphics command.
%\graphicspath{{examples/documentation/images/}{images/}} % Paths in which to look for images

\usepackage{setspace} % change spacing within the document
\usepackage{framed} % framed or shaded regions that can break across pages
\usepackage{enumitem} % control layout of itemize, enumerate, description
\usepackage{caption} % customising captions in floating environments
\usepackage{subcaption} % support for sub-captions
%\usepackage{schemata} % print topical diagrams
\usepackage[all]{xy} % A package for typesetting a variety of graphs and diagrams with TeX

\usepackage{amsmath} % mathematical typesetting
\usepackage{amssymb} % mathematical symbols
\usepackage{mathrsfs} % \mathscr script font that is more elaborate than the ``calligraphic'' font obtained by \mathcal
\usepackage{extarrows} % extra Arrows beyond those provided in amsmath
\usepackage{mathtools} % enhance the appearance of documents containing a lot of mathematics
\usepackage{cancel} % a package to draw diagonal lines (“cancelling” a term) and arrows with limits (cancelling a term “to a value”) through parts of maths formulae.
\usepackage[official]{eurosym} % euro symbol
\usepackage{halloweenmath} % Scary and creepy math symbols
%\usepackage{svrsymbols} % A font with symbols for use in physics texts
%\usepackage{mhchem} % typeset chemical formulae/equations and Risk and Safety phrases
%\usepackage{bohr} % simple atom representation according to the Bohr model
%\usepackage{marvosym} % symbols for structural engineering; symbols for steel cross-sections; astronomy signs (sun, moon, planets); the 12 signs of the zodiac; scissor symbols; CE sign and others.

\usepackage{comment}
\usepackage[backend=bibtex,style=alphabetic,natbib=true]{biblatex} % Use the bibtex backend with the authoryear citation style (which resembles APA)
\usepackage[autostyle=true]{csquotes} % Required to generate language-dependent quotes in the bibliography
%\addbibresource{sample.bib}  %The filename of the bibliography


\geometry{lmargin=2.5cm,rmargin=2.5cm,tmargin=3cm,bmargin=2.5cm}

\setlength{\parindent}{0pt} % remove indent

%\sloppy % rules to fix some line breaking
%\fuzzy % default LaTeX rules
%\date{}


\begin{document}

% Uncomment next three lines if you want default values



%-----------------------------------
%           TITLE PAGE
%-----------------------------------

\begin{titlepage}
 \newcommand{\HRule}{\rule{\linewidth}{0.5mm}}

	\center

	%------------------------------------------------
	%	Headings
	%------------------------------------------------

    \textsc{\LARGE \bf\textcolor{red}{Church of Emacs}}\\[1.5cm] % Main heading such as the name of your university/college

    \textsc{\Large \textcolor{red}{Pagan Estudies}}\\[0.5cm] % Major heading such as course name

	\textsc{\large Vi IMproved}\\[0.5cm] % Minor heading such as course title

	%------------------------------------------------
	%	Title
	%------------------------------------------------

	\HRule\\[0.4cm]

    {\huge\bfseries $\mathwitch$ Number VI-VI-VI $\reversemathwitch$}\\[0.4cm] % Title of your document

	\HRule\\[1.5cm]

	%------------------------------------------------
	%	Author(s)
	%------------------------------------------------

	\begin{minipage}{0.4\textwidth}
		\begin{flushleft}
			\large
            \textit{\textcolor{red}{Autor}}\\
			cyma % Your name
		\end{flushleft}
	\end{minipage}
	~
	\begin{minipage}{0.4\textwidth}
		\begin{flushright}
			\large
            \textit{\textcolor{red}{Supervisor}}\\
			Flying Spaghetti Monster% Supervisor's name
		\end{flushright}
	\end{minipage}

	% If you don't want a supervisor, uncomment the two lines below and comment the code above
	%{\large\textit{Author}}\\
	%The \textsc{Coconut} % Oil

	%------------------------------------------------
	%	Logo
	%------------------------------------------------

	%\vfill
	%\includegraphics[width=0.3\textwidth]{/path/to/picture.jpg}\\[1cm] % Include a department/university logo.

	%---------------------------------------------------------------------------


    %------------------------------------------------
	%	Date
	%------------------------------------------------

	\vfill % Position the date 3/4 down the remaining page

	{\large\today} % Date, change the \today to a set date if you want to be precise


    \vfill % Push the date up 1/4 of the remaining page

\end{titlepage}



%-------------------------------------------------------------------------------
%	     CONTENT - CHAPTERS
%-------------------------------------------------------------------------------

\pagenumbering{roman}


\pagestyle{fancy}


\renewcommand{\sectionmark}[1]{\markright{\thesection.\ #1}} %header conf. to show section

\renewcommand{\footrulewidth}{0pt}
\renewcommand{\headrulewidth}{0.4pt}



\fancyhead{}

\fancyfoot{}
\fancyfoot[RE]{\thepage}
\fancyfoot[LO]{\thepage}

\renewcommand{\contentsname}{Contenido}

\doublespacing
    \tableofcontents % Prints the main table of contents
        \addtocontents{toc}{~\hfill\textbf{\small{Página}}\par} % add pages to toc
\singlespacing

%-------------------------------------------------------------------------------
%                   PRE-SECTIONS
%-------------------------------------------------------------------------------

\newpage

% In case you need to add more lists, do it here

\renewcommand{\listfigurename}{Lista de figuras}
\renewcommand{\listtablename}{Lista de cuadros}

\listoffigures % Prints the list of figures move
    \addcontentsline{toc}{section}{\listfigurename}

\listoftables % Prints the list of tables
    \addcontentsline{toc}{section}{\listtablename}

\newpage

\section*{Introducción}
    \addcontentsline{toc}{section}{Introducción}


%---------------------------------------------------------------------------------------
%               BEGIN SECTIONS AREA
%---------------------------------------------------------------------------------------

\newpage

\pagenumbering{arabic}

\fancyhead[RE,LO]{\textsl{\rightmark}} % Show the section/subsection in header

% Comment next line in case you need to reset equation counter in each section
%\counterwithin*{equation}{section}

\section{Exit VIM}



\newpage
\section{Celebrate Pirate's Day}


\newpage

%---------------
%   REFERENCES
%---------------
%Add heading=bibintoc or subbibintoc to add in table of contents.

%\printbibliography[type=article,title={Articles only}]
%\printbibliography[type=book,title={Books only}]
%\printbibliography[keyword={physics},title={Physics-related only}]

%\printbibliography[heading=bibintoc,type=online,title={Webgrafía}]



\end{document}


